\section{Courage and Fearlessness}

\begin{quotex}
The would-be initiate must bring with him a certain measure of courage and fearlessness. He must positively go out of his way to find opportunities for developing these virtues. His training should provide for their systematic cultivation. In this respect, life itself is a good school --- possibly the best school. The student must learn to look danger calmly in the face and try to overcome difficulties unswervingly. For instance, when in some peril, he must swiftly come to the conviction that fear is of no possible use: I must not feel afraid; I must only think of what is to be done. And he must improve to the extent of feeling, upon occasions which formerly inspired him with fear, that to be frightened, to be disheartened, are things that are out of the question as far as his own inmost self is concerned. By self-discipline in this direction, quite definite qualities are developed which are necessary for initiation into the higher mysteries.
\flright{\textsc{Rudolph Steiner}, \textit{Knowledge of the Higher Worlds}}
\end{quotex}

In \textit{Sintesi di dottrina della razza} (Synthesis of Racial Doctrine), Evola included several pages of photographs illustrating various aspects of Aryan men. He included not one, but two, photos of Rudolph Steiner with this caption:

\begin{quotex}
Among the race of the nordic-dinaric body type, we see the solar element differentiating itself, predominantly in the ``ascetic'' sense ... The physical traits give the sense of a condition of illumination and a power of spiritual penetration.
\end{quotex}

\flrightit{Posted on 2010-02-27 by Cologero}

\begin{center}* * *\end{center}

\begin{footnotesize}
\begin{sffamily}
\texttt{Tom Blanchard on 2014-08-21 at 10:56 said:}

``I must not fear. Fear is the mind-killer. Fear is the little-death that brings total obliteration. I will face my fear. I will permit it to pass over me and through me. And when it has gone past I will turn the inner eye to see its path. Where the fear has gone there will be nothing. Only I will remain.''

Frank Herbert, Dune

\hfill

\texttt{\_ on 2015-08-17 at 02:27 said:}

Overcoming irrational fears would be a good place to start for those like me who are too careful. Think of the fear of sharks, totally irrational but very gripping on the minds of people!

\hfill

\texttt{Cologero on 2015-08-17 at 11:59 said:}

@swimming with sharks:

You should check this out:

\textit{Qualifications for Initiation}\footnote{\url{http://www.gornahoor.net/?p=1161}}.

Perhaps it is more than just symbolic.

\hfill

\end{sffamily}
\end{footnotesize}

